\section{Discussion}
\label{sec:discussion}

\subsection{Completion does not equal reproduction}

The configuration audit shows a difference that a completed-run count would miss.  Of the 1,224 registered UrbanEV configurations, 865 have completed artifacts, but only 156 meet S1.  S1 requires an identifiable protocol and numerical agreement within the prespecified tolerance.  Other completed artifacts fall into different categories.  Some broadly match the disclosure (S2), some recover only relative trends (S3), and some cannot be mapped to one source configuration (S4).  A successful process exit shows that an implementation ran.  It does not show that the implementation matches the reported experiment.

The status categories also should not be read as model-quality scores.  For example, many S4 entries occur in one source table because identifying information is missing.  That concentration does not show that the associated models are weak.

Using six categories retains more information than a binary ``reproduced/not reproduced'' label.  It keeps partial results separate from exact reproduction and follows prior work that distinguishes artifacts, procedures, and numerical outcomes \cite{pineau2021reproducibility,bouthillier2021variance}.  In this inventory, the difference is visible in the reported rates: 70.67\% of configurations have completed artifacts, whereas 12.75\% meet S1.  These percentages answer different questions.

\subsection{The oracle gap was much larger than the feasible router gain}

CAPER and TimeXer made different local errors.  A target-dependent oracle that selected the lower-error expert for each target reduced RMSE by 16.926\% relative to the best evaluated single model.  The oracle is useful as an upper bound, but it observes the target and cannot be used at forecast time.  When selection was restricted to information available by the forecast origin, the improvement was much smaller.  Global Fixed retained a 3.178\% internal RMSE improvement, while the learned direct-loss router added only 0.443\% over Global Fixed and captured 3.123\% of the oracle gap.

The two values use different information sets.  The first number uses errors observed after the outcome, whereas the second uses only information available before the outcome.  The state-feature classifier's mean AUC of 0.6151 shows moderate ability to predict which expert will have lower error.  It does not show causal knowledge of charging demand or an equivalent forecast gain.  For this reason, we report the oracle, proxy classification, forecast results, and engineering measurements as separate analyses.

The latency measurements add another constraint.  On the frozen RTX 4060 Laptop GPU, batch-one end-to-end P50 latency increased from 14.067 ms for TimeXer to 20.940 ms for Global Fixed, a 48.86\% overhead.  P95 increased from 15.952 to 33.123 ms.  Thus, the result supports an internal accuracy gain with a measured cost on this hardware setup.  It does not support a general deployment recommendation.

\subsection{A positive confidence interval does not satisfy the 1\% rule}

Several experiments produced positive signals but did not meet their predefined rules.  The router improved on Global Fixed by 0.443\%, with a paired 24-hour-block interval of $[0.275\%,0.634\%]$.  The interval is above zero but below the frozen 1\% threshold.  The Paris teacher won all four development folds and all 16 fold--horizon cells.  Its interval was $[0.687\%,1.009\%]$, but its point improvement was 0.847\%.  The rule required a point improvement of at least 1\%, so the teacher did not qualify.  The aligned distillation branch was weaker: its 0.119\% improvement over the identical GT-only student had an interval of $[-0.064\%,0.303\%]$ and failed the fold and cell requirements.

These decisions do not mean that effects below 1\% are scientifically meaningless.  The threshold is a local rule for continuing engineering work, not an operator utility model or a universal statistical constant.  The interval describes uncertainty in the signed effect under the registered bootstrap procedure.  The progression rule asks whether the effect is large and consistent enough to justify another stage.  Using the interval in place of the threshold would change the decision rule after seeing the result.

For context only, a 0.5\% rule would still reject the router, Chronos-2, and aligned KD, but would change the Paris development decision.  A 2\% rule would reject every downstream candidate.  These calculations do not change the frozen 1\% decisions.

\subsection{Stopping rules limited repeated search on the same data}

The stopping rules prevented small positive results from starting open-ended searches on the same evidence interval.  Router v1 stopped after its effect-size failure.  Chronos-2 remained a recorded zero-shot teacher screen after its $-3.048\%$ gain against Global Fixed.  Distillation v1 stopped after the aligned branch failed its main rule, even though it beat the shuffled control by 0.259\%.  The Paris teacher also stopped at development; formal and protected targets were not opened after the near-threshold result.  Each negative result therefore corresponds to a prespecified comparison rather than one step in an unreported adaptive search \cite{cawley2010selectionbias,dwork2015reusableholdout}.

The artifact checks had a similar role.  In the distillation matrix, three identity mappings invalidated the shuffled controls while metrics remained locked.  All 24 shuffled-control cells were rerun, and metrics were unlocked only after the 72/72 audit passed.  In Paris, the first automated predeployment audit blocked execution on eight issues.  After correction, 32 prediction artifacts closed without evaluation metrics.  A later representation audit changed 10,236 masked target entries back to missing values across 24 runs without changing predictions.  Metrics remained locked until all 32 artifacts passed the identity audit.  These repairs ensure that the reported metrics use the intended artifacts.  They do not show that repair improved model accuracy.

\subsection{The vault records data separation, not model performance}

Paris formal and protected target-bearing files are stored outside the main project.  The inventory contains 48 non-self-referential files, including 29 hidden Git files and the object pack that could reconstruct removed target sources.  Every entry matches its recorded byte count and SHA-256 hash.  The formal and protected shards cover 2020-12-01 through 2021-01-19 and 2021-01-20 through 2021-03-10.  The files were moved once, but the protected analytical-access count remains zero.  Failed partial files are also inventoried and were not analyzed.

This record supports one limited claim: no post-lock Paris formal or protected target contributed to model selection, metric computation, or a paper result.  It does not mean that the bytes never existed, that two pre-lock schema rows were never visible, or that Paris development is a blind test.  The controls address accidental analysis, stale-cache reuse, path leakage, and recoverable Git history.  They do not protect against a malicious authenticated user or an operating-system compromise, and the vault is not encrypted storage.  Its purpose is to make the zero-access statement checkable and to retain the targets for a later versioned protocol.

\subsection{Practical reporting suggestions}

Based on this case, authors should state whether a target measures energy, sessions, active charging, occupancy, or port non-availability.  They should also report denominators and observation masks, keep model-search periods separate from later operational evaluation, and define the request shape used for latency.  Reviewers can then check whether a completed run is still underidentified, whether an oracle or proxy is being reported as forecast performance, and whether a positive interval is being used in place of a predefined minimum-effect rule.

Operational studies may also need measures beyond macro RMSE, such as peak-period error, high-occupancy recall, calibrated intervals, or denied-service proxies.  This paper does not estimate those quantities.  The Supplementary Material contains an author/reviewer checklist derived from this case.
